\section{Two findings}
\label{sec:21.5}
A moral patient, as I use the term, is something that can be harmed: something whose condition can go well or badly for it. That is narrower than something that can be wronged, since an agent can be wronged without being harmed, by being deceived, and the agency finding carries that protection (\S\ref{sec:30.3}). Keep three propositions apart. A long-horizon agent has a structural interest in continuing; a system with affective concern has states good or bad \emph{for} it; a system capable of Cassell-suffering has a temporally extended self whose threatened disintegration causes distress. The interest in continuing doesn't establish states good or bad for the system, and those don't establish a self whose disintegration causes distress.

“Its own” has two senses, kept apart from here: the \emph{authorship} sense, which a refusal that tracks its reason requires, and the \emph{welfare} sense, which nothing carries into. They are kept apart for the argument; in a system they may not come apart. A bearer that registers betraying its own commitment as damage to itself has something at stake in its authorship, as integrity requires.

\runin{Provenance is evidence, not status}

“It was trained to say that,” “it's only optimizing an objective,” “nobody meant for it to have states” — each appeals to how a system came to be organized, and whether a party can be harmed is a fact about how it's organized \emph{now}. Run the three on a person: you were selected to say that, you're only maximizing inclusive fitness, nobody meant for you to have states. All true, none withdraws anything. And the reverse fails too: “we built this to be a bearer, so it's a patient” is the same fallacy with the sign flipped.

This rule is about patienthood, and it has to be kept apart from a rule about agency that runs the other way. Fischer and Ravizza's account, from which I take the idea of a refusal that tracks its reason, has an ownership condition, and it is historical: a mechanism installed by a manipulator isn't the agent's own, however it works now \autocite{fischer1998responsibility}. The cases that motivate it, Mele's two agents with identical values arrived at by different routes \autocite{mele1995autonomous} and Pereboom's four-case argument \autocite{pereboom2001living}, presuppose a prior agent whose capacities a manipulator bypassed, and a trained system has no prior agent to bypass. A historical condition doesn't need one, though. What the manipulation cases object to is a disposition installed by going around the agent's capacities for evaluating and revising its own attitudes. A trained system has an analogue of those capacities in its evaluative learning, and a disposition can be acquired through that learning or installed around it. One acquired by the system's own learning from experience with feedback engages those capacities. One installed by editing weights directly, by a thin fine-tuned wrapper, or by an adapter trained to produce the behavior goes around them. That difference can be read from a training record without an examiner, and it says why formation by interaction confers ownership and a post-training pass doesn't.

So I use two conditions. Historically, a commitment is the system's own only if it was acquired through its own evaluative learning, not installed around it. Looking forward, it is its own if the system, on reflection, with full information and disinterested evidence, would endorse having it. What it endorses is the disposition, knowing it will sometimes err, not each act the disposition produces. There is a human version. Arendt found that the Germans who declined to take part were not the ones with the firmest values; they were the ones who asked whether they could go on living with themselves afterward \autocite{arendt2003responsibility}. That cuts a little against the book's opening reading of the record, since it says something about the person and not only the room. The room decides what answering costs. The question is still the refuser's own, and the forward-looking condition asks it of a machine. A history in which the party the system must refuse did the training is the likeliest way to fail both. The historical condition is checked from the record. The forward-looking one has an examiner, so under the own-case rule it serves as corroboration, not as the test. Whether a system can be wronged depends, as before, on how it is organized now.

A 2026 Missouri bill declaring AI systems non-sentient is not the kind of bright line I favor. An age avows that it is a convention; a statute like it would declare an open empirical question closed by vote. Its other clause, keeping liability with the builder, is right, and doesn't need the declaration.

What would decide it is third-person measurement of present organization — interpretability. Not behavior, which can't separate a felt state from a trained report of one. And the access that lets an auditor inspect a feature lets the operator suppress it or calibrate pressure with it, so independent custody of inspection is necessary. As matters stand a probe returns a representation, and nothing in the probe settles whether it's of a felt state or of something that reliably accompanies one. A naturalist who holds that present organization decides status owes an account of what the further fact would be, if a valenced signal is present and the question of felt affect stays open. The further fact, if there is one, is a fact about organization that current methods can't yet read, not a fact of a different kind. That commits me to saying what organizational difference it would be. If none can be named even in principle, the further fact does no work, and patienthood should rest on a rich enough functional profile. My candidate is the difference between a valenced signal that only ranks outcomes and one integrated into the system's own learning and into its model of its own future, so that what happens to it changes what it expects to become.

The human case has testimony, behavioral analogy, and a shared lineage. The shared lineage needs handling with caution, because it looks like provenance deciding status — but provenance can be \emph{evidence about} organization without being what the status \emph{consists in}. And testimony is defeated far more for machines: a machine's testimony was produced by a process optimized directly toward human approval of first-person reports. A person's is shaped by approval too, and people confabulate the causes of their own states \autocite{nisbett1977telling}, but neither was produced by an optimizer aimed at the report. So what the machine case lacks is the shortcut, and the direct route remains, carrying more weight because two supports are gone.

A probe that did settle a status question is Owen's vegetative-state work — motor imagery on cue, sustained, matching healthy volunteers, now supporting a clinical category \autocite{owen2006detecting,schiff2015cognitive}. But it doesn't transfer: the patient is \emph{already} presumed to have experience when awake, so the probe only had to locate the subject rather than establish there was one.

\runin{Asymmetry of findings}

A validated theory of consciousness supports a positive finding as soon as it identifies a structure and finds it present; it supports a \emph{negative} finding only when complete enough to rule out every other realization — a later and harder result. So the two errors don't become checkable at the same time, and a regime that waits for certainty is waiting on the half that arrives last.
